\section{Discussion}
\label{sec:discussion}

\subsection{Geometry-Dependence of Algorithm Defaults}
\label{sec:discussion:geometry}

The published Table~1 hyperparameters of \mcpilot{}~\cite{turcato2025mcpilot} are
calibrated to a Franka Panda arm at a fixed release height with target range
$\approx 1.0$\,m.
Our study quantifies how several defaults degrade silently when geometry changes.

The most important default is the RBF lengthscale $\ell_s$.
Config~A has $\ell_s = R$ by coincidence; once $R$ deviates from 1.0\,m the constant
is no longer appropriate.
The rule $\ell_s \approx 0.15 \times R$ (Section~\ref{sec:convergence:lengthscale})
transfers across all five tested heights without per-height tuning beyond verifying
that $S < 0.1$ at the domain extremes.

The cost lengthscale $\ell_c = 0.1$\,m also requires attention.
With random exploration, early landing errors span 0.3--1.5\,m, saturating the cost at
$c \approx 1.0$ and driving gradients to zero.
Widening to $\ell_c = 0.5$\,m keeps gradients alive across the full exploration error
range; further tightening after early convergence is optional.

\subsection{Two Faces of Disturbance Handling}
\label{sec:discussion:disturbances}

Studies~3 and~4 together reveal a fundamental dichotomy in how the GP handles disturbances.

\paragraph{Zero-mean noise cannot be conditioned out.}
Theorem~\ref{thm:zeromean} establishes that symmetric zero-mean noise does not shift
the optimal mean action.
The GP will faithfully represent the increased aleatoric variance in its posterior,
but no amount of noise-aware training improves mean performance.
This is why pure Gaussian additive noise is unlearnable: the aware and naive policies
reach the same optimum.
Only \emph{biased} disturbances --- velocity slip, timing jitter, salt-and-pepper
spikes --- shift the mean landing and are learnable.
The original \mcpilot{} hardware model is a timing jitter with $a > 0$, which is
biased, explaining why it improves performance on real hardware.

\paragraph{Biased disturbances can be conditioned out, but only when the trial budget
supports the extra GP dimensions.}
The wind study shows that explicit wind conditioning fails within the 15-trial budget.
The blind GP treats constant wind as an implicit dynamics bias and absorbs it without
any state augmentation.
For slowly varying wind this is strictly better than explicit conditioning in the
low-data regime.
The transition point is approximately $\sim 40$ trials for a 4-D policy space.

These two findings together suggest a practical hierarchy for disturbance handling:
(1) if the disturbance is zero-mean symmetric, do not add it to the state (useless);
(2) if biased and constant, let the GP absorb it implicitly (free);
(3) if biased and time-varying with a sufficient trial budget, augment the state.

\subsection{The Decoupled Architecture as Portability Mechanism}
\label{sec:discussion:decoupled}

Setting ball velocity via \texttt{resetBaseVelocity} at release decouples arm tracking
quality from GP training quality.
Study~2 demonstrates the consequence: 100\% commanded-velocity hit rates on all three
arms without retraining, even though the KUKA iiwa7 EE velocity tracking error is
$\approx 7\%$.
The 33\% achieved hit rate on all arms reflects arm tracking error that would be
present on any real robot using a standard position controller.

In a physical deployment, this decoupling corresponds to a gripper that can measure the
EE velocity at the release instant (via encoder forward kinematics) and set ball velocity
to the measured value.
In that case, the delivered velocity matches the commanded velocity, and the RL
algorithm's data-efficiency guarantees hold regardless of trajectory-tracking quality.

\subsection{Limitations}
\label{sec:discussion:limitations}

\paragraph{Simulation only.}
All experiments are conducted in simulation.
Hardware validation on a physical manipulator is needed to confirm that learned policies
transfer and that the 5--10\% PyBullet--NumPy drag discrepancy does not compound with
real-world effects.

\paragraph{Multi-seed statistics.}
Most results are reported at seed=1.
Config~B (random vs.\ stratified) was evaluated at seed=2 for robustness verification;
the seed=1 vs.\ seed=2 comparison (0/10 vs.\ 10/10) illustrates that single-seed
results should be read as existence proofs rather than expected-case performance.
Conference-level statistical claims require $\geq$5 seeds.

\paragraph{Fixed launch angle.}
The elevation angle $\alpha = 35^\circ$ is fixed across all experiments.
Joint optimisation of $\alpha$ and release speed would increase the reachable target
volume, particularly for near-field targets where the current angle is sub-optimal.

\paragraph{Wind trial budget.}
The 15-trial budget is insufficient for the 4-D wind-aware policy space.
Our analysis (Section~\ref{sec:study4:dimscaling}) suggests $\sim$40 trials are needed.
Running full-budget wind-aware experiments remains future work.

\paragraph{Vision quantitative evaluation.}
The vision pipeline (Study~5) is described and implemented but not quantitatively
evaluated across a randomised target distribution with multi-seed statistics.
Measuring detection-to-throw accuracy across 50+ targets remains future work.

\paragraph{Particle freeze approximation.}
Frozen particles introduce $\approx 2$--6\,cm of systematic over-estimation in the
landing distance (Appendix~\ref{app:freeze}).
This biases cost estimates slightly pessimistically but does not prevent convergence.
